\documentclass[11pt,a4paper]{article}

\usepackage[utf8]{inputenc}
\usepackage[spanish]{babel}
\usepackage{geometry}
\usepackage{hyperref}
\usepackage{enumitem}
\usepackage{titlesec}
\usepackage{booktabs}
\usepackage{array}
\usepackage{xcolor}
\usepackage{fancyhdr}
\usepackage{parskip}
\usepackage{graphicx}
\usepackage{multicol}

\geometry{margin=2.5cm}

\definecolor{primary}{RGB}{28, 66, 232}
\definecolor{secondary}{RGB}{232, 194, 71}
\definecolor{dark}{RGB}{30, 58, 95}
\definecolor{gray}{RGB}{120, 120, 120}
\definecolor{green}{RGB}{34, 139, 87}

\pagestyle{fancy}
\fancyhf{}
\rhead{\textcolor{gray}{\small MexGo — Fase 2}}
\lhead{\textcolor{gray}{\small Genius Arena · Fundación Coppel 2026}}
\rfoot{\textcolor{gray}{\small \thepage}}
\renewcommand{\headrulewidth}{0.3pt}

\titleformat{\section}{\large\bfseries\color{primary}}{}{0em}{}[\titlerule]
\titleformat{\subsection}{\normalsize\bfseries\color{dark}}{}{0em}{}
\titleformat{\subsubsection}{\normalsize\bfseries\color{gray}}{}{0em}{}

\newcommand{\criterio}[2]{
  \vspace{0.3cm}
  \noindent\colorbox{primary!10}{\parbox{\dimexpr\linewidth-2\fboxsep}{%
    \textbf{\textcolor{primary}{#1}}\\[2pt]
    \small\textcolor{dark}{#2}
  }}
  \vspace{0.3cm}
}

\newcommand{\feature}[1]{\textcolor{green}{\textbf{✓}} #1}
\newcommand{\tech}[1]{\texttt{\textcolor{primary}{#1}}}

\hypersetup{
  colorlinks=true, linkcolor=primary, urlcolor=secondary,
  pdftitle={MexGo FASE 2}, pdfauthor={Los Mossitos}
}

% ===========================================================================
\begin{document}
% ===========================================================================

{\centering
  {\Large\bfseries\color{primary} MexGo}\par
  \vspace{0.2cm}
  {\normalsize\color{dark} Plataforma multilingüe de turismo inclusivo para el Mundial 2026}\par
  \vspace{0.2cm}
  {\small\color{gray} Los Mossitos · Genius Arena · Fundación Coppel 2026 · \today}\par
  \vspace{0.5cm}
  \hrule
  \vspace{0.8cm}
}

\tableofcontents
\vspace{0.8cm}
\hrule
\vspace{0.5cm}

% ===========================================================================
\section{Resumen ejecutivo}
% ===========================================================================

MexGo es un puente digital diseñado para que la derrama económica del Mundial FIFA 2026
llegue a quienes suelen ser invisibles: los micro y pequeños negocios locales.
Nuestra plataforma conecta a turistas internacionales con la autenticidad de México,
mientras integra a los negocios en el ecosistema de capacitación de \textbf{Coppel Emprende}.

Utilizamos \textbf{IA (Gemini)} no solo para traducir o guiar, sino como el motor de equidad
que decide qué negocios recomendar basándose en su capacidad, accesibilidad y el perfil
cultural del turista. El resultado es un prototipo funcional que transforma a una fonda
tradicional o un taller de artesanías en un destino global, eliminando las barreras de
idioma y tecnología sin fricción para el usuario.

% ===========================================================================
\section{Problemática y contexto}
% ===========================================================================

En 2026, el Mundial traerá millones de dólares a México, pero el gran problema es el contraste:
un turista inglés, como Thomas, pasando de largo frente a la fonda de Doña Carmen para ir a
una franquicia. ¿Por qué? Porque no hablan el mismo idioma y ella es invisible en internet.

A pesar del esfuerzo de micronegocios, 8 de cada 10 mueren antes de los dos años por
falta de clientes y de conocimiento digital. Esta desigualdad económica hace que las
grandes cadenas acaparen los beneficios mientras el pequeño comercio solo se queda
mirando. La barrera no es la calidad del servicio, sino la \textbf{invisibilidad digital}
y la \textbf{brecha de capacitación}.

% ===========================================================================
\section{Solución: MexGo}
% ===========================================================================

MexGo es el puente digital que conecta la cartera del turista directamente con
el micronegocio local. Nuestra solución no es solo un buscador; es un ecosistema
que profesionaliza al negocio mientras atrae al cliente.

\subsection{Para el turista: Experiencia auténtica y accesible}

\begin{enumerate}[leftmargin=*, itemsep=3pt]
  \item \textbf{Acceso sin fricción}: PWA rápida, multilingüe (ES, EN, FR) vía QR.
  \item \textbf{Recomendaciones por IA}: Gemini recibe el perfil del turista (país,
        grupo, accesibilidad) y \textbf{decide dinámicamente} qué negocios sugerir.
        La equidad se inyecta en el prompt: la IA prioriza negocios con menor
        afluencia actual o con certificaciones específicas (LSM).
  \item \textbf{Asistente de viaje}: Chat 24/7 que traduce menús, explica
        costumbres locales y genera itinerarios inteligentes integrando
        el Mundial.
\end{enumerate}

\subsection{Para el negocio: Visibilidad y Crecimiento}

\begin{enumerate}[leftmargin=*, itemsep=3pt]
  \item \textbf{Registro simplificado}: Onboarding rápido para digitalizar
        la oferta básica del negocio.
  \item \textbf{Capacitación Coppel Emprende}: Acceso directo a módulos de
        formación para dejar de ser invisibles. Incluye certificación en
        \textbf{Lengua de Señas Mexicana (LSM)} para abrir líneas de ingresos
        que la competencia ignora.
  \item \textbf{Visibilidad asistida}: La IA traduce automáticamente las
        descripciones y facilita la comunicación con clientes de cualquier
        nacionalidad, eliminando la barrera del idioma.
\end{enumerate}

% ===========================================================================
\section{Criterios de evaluación}
% ===========================================================================

\subsection{Alto impacto — relación con el track}

\criterio{Track: negocios turísticos locales en contexto del Mundial 2026}
{MexGo ataca directamente la problemática: micronegocios sin visibilidad
internacional y turistas sin guía cultural adaptada.}

\begin{itemize}[leftmargin=*, itemsep=2pt]
  \item El torneo trae turistas de decenas de países con idiomas y culturas
        distintas; MexGo es el único punto de entrada adaptado por perfil cultural.
  \item La \textbf{equidad impulsada por IA} prioriza negocios locales sobre grandes cadenas,
        garantizando que el flujo económico llegue al nicho objetivo de Fundación Coppel.
  \item Accesibilidad como primer ciudadano: el cuestionario captura necesidades
        (discapacidad auditiva, visual, motriz) y filtra recomendaciones en consecuencia
        — un diferenciador que ninguna plataforma comercial ofrece de base.
  \item Información práctica inmediata: el turista sale del cuestionario con
        tarjetas listas, no con una lista de 200 resultados sin contexto.
\end{itemize}

\subsection{Innovación y creatividad}

\criterio{¿Qué lo hace diferente?}
{Combinación de perfil cultural + accesibilidad + IA multilingüe aplicada
a negocios locales en un flujo sin fricción.}

\begin{itemize}[leftmargin=*, itemsep=2pt]
  \item \textbf{Equidad por Diseño e Inteligencia}: el sistema no recomienda por popularidad,
        sino por equidad. La IA analiza el perfil del visitante y el estado de los negocios
        (saturación, certificaciones Ola México) para equilibrar la demanda. Un turista
        japonés recibe sugerencias distintas a un alemán, asegurando que la ganancia
        sea justa y distribuida.
  \item \textbf{PWA multilingüe con i18n dinámica}: la interfaz, el chat y las
        recomendaciones cambian de idioma en tiempo real. No hay tres apps —
        hay una que se adapta.
  \item \textbf{IA embebida como motor de decisión}: Gemini no es solo un chat; es el
        motor que decide qué negocios mostrar basándose en las reglas de negocio
        inyectadas en su prompt de sistema.
  \item \textbf{QR de entrada}: el turista llega desde un volante, un letrero
        en el estadio o un stand de Fundación Coppel. Sin descargar nada.
\end{itemize}

\subsection{Escalabilidad}

\criterio{De hackathon a producto nacional y regional}
{La arquitectura fue diseñada explícitamente para escalar post-hackathon.}

\begin{itemize}[leftmargin=*, itemsep=2pt]
  \item \textbf{Modelo de datos extensible}: Supabase + migraciones versionadas
        (\tech{supabase/migrations/}). Agregar ciudades, tipos de negocio o
        idiomas es una migración y un archivo de traducciones.
  \item \textbf{Algoritmo desacoplado}: \tech{lib/equity.ts} es una función
        pura — se puede mover a una Supabase Edge Function o a un microservicio
        sin tocar el frontend.
  \item \textbf{Multiidioma por diseño}: \tech{next-intl} con archivos JSON
        independientes. Agregar portugués para Brasil o alemán para la delegación
        europea es un archivo de 300 líneas.
  \item \textbf{Sin servidor propio}: desplegado en Vercel (serverless). Escala
        automáticamente con el tráfico del mundial sin aprovisionar
        infraestructura.
  \item \textbf{Módulo de negocios replicable}: el mismo dashboard puede
        habilitarse para otros eventos deportivos, culturales o de turismo
        regional en México.
\end{itemize}

\textbf{Ruta de crecimiento post-hackathon:}
\begin{enumerate}[leftmargin=*, itemsep=1pt]
  \item Piloto con negocios reales en una sede (ej. Guadalajara o Monterrey).
  \item Integración con Google Maps Platform para enriquecer datos de negocio.
  \item Módulo de pagos (Conekta / Stripe) para reservas directas.
  \item App nativa (React Native) si la demanda lo justifica.
  \item Extensión a eventos Fundación Coppel fuera del mundial.
\end{enumerate}

\subsection{Viabilidad}

\criterio{Stack coherente con el problema, equipo y presupuesto}
{Todas las herramientas usadas tienen tier gratuito suficiente para el
hackathon y pricing predecible para producción.}

\begin{center}
\begin{tabular}{>{\bfseries}p{3.5cm} p{4.5cm} p{5cm}}
\toprule
Tecnología & Rol & Costo operativo \\
\midrule
Next.js 16 + Vercel    & Frontend + API Routes + deploy & Gratis (hobby) / \$20 USD/mes (pro) \\
Supabase               & Base de datos + auth + storage & Gratis hasta 500 MB / \$25 USD/mes \\
Gemini 2.0 Flash       & Chat + itinerarios + traducción & \$0.075 USD / 1M tokens \\
next-intl              & i18n (es/en/fr)                & Open source \\
MUI + Tailwind         & UI components + estilos        & Open source \\
\bottomrule
\end{tabular}
\end{center}

\textbf{Estimado mensual para 10{,}000 turistas activos:} \textasciitilde\$60--80 USD. Para el
piloto del mundial con carga real, el costo escala de forma lineal y predecible.

El equipo de 5 personas cubrió los módulos sin dependencias externas de pago
durante el desarrollo. No se requirió hardware ni licencias.

\subsection{Mérito técnico}

\criterio{Prototipo funcional, no solo maqueta}
{El prototipo presentado corre en producción con flujos completos
de turista y negocio.}

\textbf{Funcionalidades operativas en el prototipo:}

\begin{itemize}[leftmargin=*, itemsep=2pt]
  \feature{Registro e inicio de sesión completos (turista y negocio) con JWT + cookies}
  \feature{Cuestionario de 4 pasos con validación, accesibilidad y envío a API real}
  \feature{Chat con Gemini integrado con historial y contexto de perfil}
  \feature{Dashboard de negocio con módulos de aprendizaje y gestión de solicitudes}
  \feature{Sistema de badges/insignias para negocios}
  \feature{Internacionalización dinámica español / inglés / francés}
  \feature{Modo oscuro / claro con persistencia}
  \feature{Navbar auth-aware (muestra login o perfil según sesión)}
  \feature{Middleware de protección de rutas (proxy.ts)}
  \feature{34 rutas compiladas sin errores en build de producción}
  \feature{Migraciones de base de datos versionadas (4 migraciones)}
  \feature{API REST documentada (docs/API.md)}
\end{itemize}

\textbf{Stack técnico completo:}

\begin{center}
\begin{tabular}{>{\bfseries}p{3.5cm} p{10cm}}
\toprule
Capa & Tecnologías \\
\midrule
Frontend       & Next.js 16.2 (App Router), React 19, TypeScript, Tailwind CSS, MUI v6 \\
i18n           & next-intl, archivos JSON por locale (es/en/fr) \\
Backend (BFF)  & Next.js API Routes, JWT, Supabase JS Client \\
Base de datos  & Supabase (PostgreSQL) con Row Level Security \\
IA             & Google Gemini 2.0 Flash \\
Auth           & Supabase Auth + cookies HttpOnly \\
Deploy         & Vercel (serverless, CDN global) \\
\bottomrule
\end{tabular}
\end{center}

\subsection{Estructura del proyecto (File Tree)}

El proyecto sigue un patrón de \textbf{Monolito Modular} con separación clara por responsabilidades:

\begin{multicols}{2}
\begin{verbatim}
mexgo/
├── app/
│   ├── (tourist)/      # PWA Turista
│   ├── (business)/     # Panel Negocio
│   ├── api/            # BFF (IA & Data)
│   └── layout.tsx      # Global shell
├── components/
│   ├── tourist/        # UI Turista
│   ├── business/       # UI Negocio
│   └── ui/             # Design System
├── lib/                # Lógica Pura
│   ├── gemini.ts       # Orquestador IA
│   ├── equity.ts       # Lógica Equidad
│   ├── supabase.ts     # Data Access
│   └── tools/          # Tool Calling
├── types/              # Typescript Defs
├── supabase/           # Migraciones DB
└── public/             # Assets & PWA
\end{verbatim}
\end{multicols}

% ===========================================================================
\section{Flujos y Arquitectura de Datos}
% ===========================================================================

\subsection{Flujo del Turista: De la invisibilidad a la conexión}

\begin{enumerate}[leftmargin=*, itemsep=2pt]
  \item \textbf{Ingreso (QR)}: El turista accede sin registro obligatorio.
  \item \textbf{Perfilado Inteligente}: Cuestionario de 4 pasos (País, Grupo, Accesibilidad).
  \item \textbf{Consultas por IA}: El turista pregunta en su idioma. Gemini activa \texttt{recomendar\_negocios}.
  \item \textbf{Equidad Dinámica}: La IA evalúa la base de datos de negocios activos, filtra por accesibilidad y prioriza aquellos con menor saturación o mayor avance en Coppel Emprende.
  \item \textbf{Acción}: El turista guarda el lugar en su itinerario o inicia navegación.
\end{enumerate}

\subsection{Flujo del Negocio: Crecimiento y Certificación}

\begin{enumerate}[leftmargin=*, itemsep=2pt]
  \item \textbf{Onboarding}: Registro en la red MexGo.
  \item \textbf{Diagnóstico IA}: La IA analiza la zona del negocio y sugiere cursos específicos (ej. "Atención al turista europeo" o "LSM Básico").
  \item \textbf{Aprendizaje (Coppel Emprende)}: El equipo completa módulos.
  \item \textbf{Insignias}: Al validar conocimientos, el negocio obtiene insignias públicas (Sello Ola México, Accesibilidad Pro).
  \item \textbf{Impulso en IA}: A mayor capacitación, la IA otorga mayor relevancia en las recomendaciones al turista.
\end{enumerate}

\subsection{Modelo de Datos Principal (Supabase)}

Diseño relacional optimizado para escala y seguridad (RLS):

\begin{center}
\begin{tabular}{>{\bfseries}p{3.5cm} p{10cm}}
\toprule
Tabla & Propósito \\
\midrule
\texttt{users\_profile}     & Perfiles de turistas y dueños de negocio. \\
\texttt{business\_profiles}  & Datos canónicos del negocio (coords, descripción, categoría). \\
\texttt{learning\_modules}   & Catálogo de cursos integrados con Coppel Emprende. \\
\texttt{business\_badges}    & Gamificación y certificaciones obtenidas (ej. LSM). \\
\texttt{chat\_messages}      & Historial persistente para seguimiento de IA. \\
\texttt{daily\_saturation}   & Tabla agregada para que la IA sepa la afluencia en tiempo real. \\
\bottomrule
\end{tabular}
\end{center}

\begin{enumerate}[leftmargin=*, itemsep=3pt]
  \item Escanea QR / accede a URL → app detecta idioma automáticamente.
  \item Pantalla de bienvenida con botones \textit{Descubrir} y \textit{Viajes}.
  \item Si no tiene sesión → modal de login/registro (no bloquea la navegación,
        se puede cerrar).
  \item Registro en 2 pasos → cuestionario de perfil en 4 pasos.
  \item Al terminar el cuestionario → chat con IA para explorar recomendaciones.
  \item El turista puede navegar \textit{Discover} (negocios), \textit{Trips}
        (itinerarios) y \textit{Chat} desde el navbar.
\end{enumerate}

% ===========================================================================
\section{Equipo}
% ===========================================================================

\begin{center}
\begin{tabular}{>{\bfseries}p{2.5cm} p{3cm} p{8cm}}
\toprule
Integrante & Rol & Contribución principal \\
\midrule
Emiliano & Business / Pitch & Modelo de negocio, narrativa de impacto y pitch \\
Alan     & Backend         & Arquitectura de API, Supabase, Auth y Base de Datos \\
Diego    & Frontend        & Desarrollo de la PWA, layouts y flujo de usuario \\
Xavi     & UX / UI         & Diseño de interfaz, experiencia de usuario y marca \\
Fidel    & IA              & Implementación de Gemini, prompts de equidad e itinerarios \\
\bottomrule
\end{tabular}
\end{center}

% ===========================================================================
\section{Repositorio y acceso al prototipo}
% ===========================================================================

\begin{itemize}[leftmargin=*, itemsep=2pt]
  \item \textbf{Repositorio}: \url{https://github.com/AGalanAG/mexgo}
  \item \textbf{Rama principal}: \texttt{main}
  \item \textbf{Build}: \texttt{npm run build} sin errores — 34 rutas generadas
  \item \textbf{Deploy}: Vercel (URL a compartir en presentación)
  \item \textbf{Documentación}: \texttt{docs/} — arquitectura, API, flujo,
        migraciones
\end{itemize}

% ===========================================================================
\section{Próximos pasos (Fase 3)}
% ===========================================================================

\begin{itemize}[leftmargin=*, itemsep=2pt]
  \item \textbf{Integración Coppel Emprende}: Conexión directa con la API de Coppel
        Emprende para sincronizar el progreso de capacitación de los negocios.
  \item \textbf{Prueba Piloto}: Lanzamiento controlado con una red de 50 negocios en
        la zona de Santa Fe / Centro Histórico.
  \item \textbf{Certificación LSM Pro}: Módulo avanzado de Lengua de Señas Mexicana
        integrado en el chat para interacciones en tiempo real.
  \item \textbf{Escalamiento Nacional}: Adaptación de la plataforma para otras sedes
        del Mundial (Guadalajara y Monterrey) y eventos de Fundación Coppel.
  \item \textbf{Análisis de Impacto}: Dashboard para Fundación Coppel que mida
        el incremento real en las ventas de los negocios participantes.
\end{itemize}

% ===========================================================================
\end{document}
